\section*{Introduction}
The study of the relation between the combinatorial/analytic torsion of
a flat vector bundle and the Morse-Smale flow was initiated by Fried \cite{Friedconj} and
Sánchez-Morgado \cite{MorgadoMorseSmale}.
In this paper, we give a formula relating
\begin{itemize}
\item a spectral invariant: the Ray-Singer metric associated
with a flat vector bundle with a Hermitian metric on a closed Riemannian
manifold;
\item a dynamical invariant: the Milnor metric which reflects
the interactions between the fixed points and the
closed orbits of the Morse-Smale flow, and generalizes the
absolute value at zero point of the Ruelle dynamical zeta function;
\item a transgressed Euler class:
the Mathai-Quillen current.
\end{itemize}

\subsection{Background}
Let $X$ be a connected closed smooth manifold of dimension $m$. Let
$(F,\nabla^{F})$ be a complex
flat vector bundle of rank $r$ on $X$ with flat connection
$\nabla^{F}$. Let $\rho: \pi_1(X)\rightarrow \GL_r(\bC)$ be the
holonomy representation of the
fundamental group $\pi_1(X)$.
Denote by $H^\sbullet(X,F)$ the cohomology of
the sheaf of locally constant sections of $F$, and by
$\lambda=\bigotimes_{i=0}^{m}(\det H^{i}(X,F))^{(-1)^{i}}$ the
determinant line of $H^{\sbullet}(X,F)$.

{Assume that $H^{\sbullet}(X,F)=0$ and that $F$ is equipped
with a flat metric, which is equivalent to say that its holonomy
representation $\rho$ is unitary.
The Reidemeister torsion \cite{ReidemeisterTorsion,FranzTorsion,dRTorsion} is a positive real number
defined by means of of a triangulation on $X$. However, it does not depend on the
triangulation and becomes a topological invariant. It is the first
invariant that could distinguish closed manifolds such as lens spaces
which are homotopy equivalent but not homeomorphic.}

The analytic torsion was introduced by Ray and Singer
\cite{RSTorsion} as an analytic counterpart of the Reidemeister torsion.
In order to define the analytic torsion one has to choose a
Riemannian metric on $X$. The analytic torsion is a certain
weighted alternating product of regularized determinants of the Hodge
Laplacians acting on the space of differential forms with values in
$F$.

The celebrated Cheeger-Müller theorem
\cite{Ch79,Muller78} tells us that the Ray-Singer analytic torsion
coincides with the Reidemeister combinatorial torsion. Bismut-Zhang \cite{BZ92} and Müller \cite{Muller2} simultaneously
considered generalizations of this result.
Müller \cite{Muller2} extended his result to the case where $F$ is
unimodular, \ie $|{\det\,}
\rho(\gamma)|=1$ for all $\gamma\in \pi_{1}(X)$. Bismut and Zhang \cite[Th.\,0.2]{BZ92} generalized the original Cheeger-Müller
theorem to arbitrary flat vector bundles with arbitrary Hermitian
metrics.
There are also various extensions to the
equivariant case by Lott-Rothenberg \cite{LottRothengerg}, Lück \cite{LuckTorsion}, and Bismut-Zhang \cite{BZ94}, to the family case by
Bismut-Goette \cite{BG01}, and to
manifolds with boundaries by Brüning-Ma \cite{BruningMa_2012}.

Let us explain Bismut-Zhang's theorem \cite[Th.\,0.2]{BZ92} in more detail. Indeed, to
formulate their result in the case where the flat vector bundle is not
necessarily acyclic or unitarily flat, Bismut and Zhang introduced the so-called Ray-Singer metric, which is a
metric on $\lambda$ defined as the product of
the analytic torsion with an $L^2$-metric on~$\lambda$.
Also they introduced the Milnor metric on $\lambda$ which is a combinatorial
metric associated with a Morse-Smale gradient flow. It generalizes the
Reidemeister torsion to the case where $F$ is neither acyclic
nor unitarily flat. In this way,
they were able to extend the Cheeger-Müller theorem to a comparison theorem of two metrics on $\lambda$, one is analytic and the other
one is combinatorial.

The study of the relation between the combinatorial/analytic torsion and the dyna\-mical system
can be traced back to Milnor \cite{MilnorZcover}. Fried
\cite{FriedRealtorsion} showed that on hyperbolic manifolds the analytic torsion of an acyclic unitarily flat vector
bundle is equal to the value at zero point of the Ruelle dynamical zeta
function of the geodesic flow. He conjectured
\cite[p.\,66~Conj.]{Friedconj} that similar results should hold true for more general
flows. In \cite{Shfried}, following previous contributions
by
Moscovici-Stanton \cite{MStorsion}, using Bismut's orbital integral formula
\cite{B09}, the author affirmed the Fried
conjecture for geodesic flows on closed locally symmetric
manifolds. In \cite{Shen_Yu}, the authors made a further
generalization to closed locally symmetric orbifolds.

Besides the gradient flow, Morse-Smale flow is the simplest structurally stable dynamical system
which has only two types of recurrent behaviors: closed orbits and
fixed points \cite{Palis68,PalisSmale70}. Fried \cite[Th.\,3.1]{Friedconj}
proved his conjecture for the Morse-Smale flows without fixed points. When compared with Bismut-Zhang's theorem \cite[Th.\,0.2]{BZ92},
it seems natural to ask whether there is a relation between the torsion invariant (or more generally the
Ray-Singer metric for non acyclic and non unitarily flat vector bundle) and a general Morse-Smale flow which has both
fixed points and closed orbits.

This is one of the motivations of Sánchez-Morgado's work
\cite{MorgadoMorseSmale}. He showed that the
heteroclinic orbits have a non trivial contribution in the torsion
invariant, and in this way he constructed a counterexample to Fried's
conjecture on Seifert manifolds.

In this
paper, we introduce a new Milnor metric,
which indeed contains the heteroclinic contributions and
generalizes the absolute value at zero point of the Ruelle dynamical zeta function, and we
give a comparison theorem for the Milnor and Ray-Singer metrics on
$\lambda$. We believe
that in this way we give a complete answer in the affirmative to the above question.

Let us mention that there is another interpretation of the Ruelle dynamical zeta
function provided by Dang-Rivière \cite{DR_MSR}. See also
\cite{DR_MSG, DR_MSI, DR_MSII, DR_MSW} for related works.

\subsection{A new Milnor metric}
A vector field $V$ is called Morse-Smale if $V$ generates a flow
whose nonwandering set is the union of a finite set $A$ of
hyperbolic fixed points and a finite set $B$ of hyperbolic closed orbits, and if the stable and unstable manifolds of the critical
elements in $A\coprod B$ intersect transversally.

Let us take a Hermitian metric $g^{F}$ on $F$. In Section
\ref{sSF}, we construct on $\lambda$ a Milnor type metric
$\|{\cdot}\|^{\rM,2}_{\lambda, V}$ using long exact
sequences associated with a Smale
filtration of the Morse-Smale flow. Note that
the long exact
sequences encode the information about the interactions
between the critical elements in $A\coprod B$. If $V$ is a negative gradient of a Morse function,
then our Milnor metric is just the classical one as defined in
\cite[Def.\,1.9]{BZ92}, which generalizes \cite{MilnorWhiteheadTor}.

Our first result says that the Milnor metric
$\|{\cdot}\|^{\rM,2}_{\lambda, V}$ is a generalization of the absolute value at
zero point of
the Ruelle dynamical zeta function. For a closed orbit $\gamma\in B$,
let $\ell_{\gamma}\in \bR^{*}_{+}$ be its minimal period, and let
$\ind(\gamma)\in \bN$ be its index
(see \eqref{indexg}). Take $\Delta(\gamma)$ to be $1$ if
$\gamma$ is untwist and $-1$ in the contrary case (see \eqref{Delta}). The Ruelle dynamical zeta function
is defined for $s\in \bC$ by
\begin{equation*}\label{intoR}
R_{\rho}(s)=\prod_{\gamma\in
B}\det(1-\Delta(\gamma)\rho(\gamma)e^{-s
\ell_\gamma})^{(-1)^{\ind(\gamma)}}.
\end{equation*}

\begin{prop}\label{prop}
If $V$ does not have any fixed points, and if none of $\Delta(\gamma)$ is an eigenvalue of $\rho(\gamma)$, then
$H^{\sbullet}(X,F)=0$, and the norm of the canonical section
$\mathbf{1}\in \bC=\lambda$ is given by
\begin{equation*}
\|\mathbf{1}\|^{\rM}_{\lambda,V}=|R_{\rho}(0)|^{-1}.
\end{equation*}
\end{prop}

\subsection{The main result of the paper}
Let $g^{TX}$ be a metric on $TX$. Let
$\psi(TX,\nabla^{TX})$ be the Mathai-Quillen current associated with the
Levi-Civita connection $\nabla^{TX}$ (see Section \ref{sMQ}). It is a
current of degree $m-1$ defined
on the total space of the tangent bundle $TX$, which takes values in
$o(TX)$, the orientation line bundle of $TX$. Let $\|{\cdot}\|^{\RS,2}_{\lambda}$ be the Ray-Singer metric
on $\lambda$ associated with $(g^{TX},g^{F})$ (see Section \ref{sRS}).
Set $
\theta(F,g^{F})=\Tr[(g^{F})^{-1}\nabla^{F}g^{F}]\in \Omega^{1}(X)$.
Our main result is the following.

\begin{thm}\label{thm1} We have
\begin{equation}\label{eqthm1}
\log
\bigl(\sfrac{\|{\cdot}\|^{\RS,2}_{\lambda}}{\|{\cdot}\|^{\rM,2}_{\lambda,V}}\bigr)=-\int_{X}\theta(F,g^{F})(-V)^{*}\psi(TX,\nabla^{TX}).
\end{equation}
\end{thm}

If $V$ does not have any closed orbits, Theorem \ref{thm1} reduces to
\cite[Th.\,0.2]{BZ92}. Note also that if $F$ is unitarily flat, then the right-hand side of
\eqref{eqthm1} varnishes. Therefore, if $V$ does not have any fixed
points and if $F$ is unitarily flat, by Proposition \ref{prop}, our theorem corresponds to \cite[Th.\,3.1]{Friedconj}.

Our proof of Theorem \ref{thm1} is based on a result of Franks
\cite[Prop.\,5.1]{Franks1979}, who
constructed a gradient flow by destroying the
closed orbits of the Morse-Smale flow.
In Section \ref{sCMM}, we first
establish a comparison formula between our Milnor metric associated with the original Morse-Smale flow
and the classical one associated with Franks'
gradient flow. In Section \ref{S2}, to obtain Theorem \ref{thm1}, we apply Bismut-Zhang's formula \cite[Th.\,0.2]{BZ92}, which compares the Ray-Singer metric with the
Milnor metric for Franks' gradient flow.

Recall that $F$ is said to be unimodular, if its holonomy representation
$\rho$ is unimodular, \ie $|{\det\,} \rho(\gamma)|= 1$ for all
$\gamma\in \pi_{1}(X)$. This is
equivalent to the fact that there is a Hermitian metric $g^{F}$
such that $\theta(F,g^{F})=0.$
By Theorem \ref{thm1}, we get
\begin{cor}
If $(F,g^{F})$ is unimodular, then
\begin{equation*}
\|{\cdot}\|^{\RS,2}_{\lambda}=\|{\cdot}\|^{\rM,2}_{\lambda,V}.
\end{equation*}
\end{cor}

\subsection{Organization of the paper}
In Section \ref{Sp1}, we
introduce some conventions on the determinant line, the cohomology
of a circle, and also a long exact sequence associated with three manifolds $Y_{1}\subset Y_{2}\subset Y_{3}$. In Section \ref{S1M}, we recall some background on Morse-Smale
flows. We also introduce the Milnor type metric, and we show Proposition
\ref{prop}. In
Section \ref{S2}, we recall the constructions of Mathai-Quillen
current and Ray-Singer metric. We show our main result.
We use the convention $\bN=\{0,1,2,\ldots\}$ and $\bR_{+}^{*}=(0,\infty)$.

\subsubsection*{Acknowledgements}
We are indebted to Xiaolong Han and Xiaonan Ma for reading a preliminary version of this paper and for useful suggestions. S.S. would like to thank Nguyen Viet Dang and Gabriel Rivière for
fruitful discussions on Morse-Smale flows.

\section{Preliminaries} \label{Sp1}
This section is organized as follows. In Section \ref{sDL}, we
introduce our convention on the determinant line. In Section
\ref{sCS1}, we give a metric on the determinant line of the
cohomology of $\bbS^{1}$. This is our model case near the
closed orbits of a flow. In~Section \ref{Fp}, we explain a long exact
sequence associated with a triple of manifolds $Y_{1}\subset
Y_{2}\subset Y_{3}$.

\subsection{The determinant line}\label{sDL}
Let $W$ be a complex finite dimensional vector space. We denote by
$W^*$ the dual space. If $\dim W=1$, we write $W^{-1}=W^{*}$. Take
$\bw\sbullet(W)=\bigoplus_{j=0}^{\dim W}\bw{j}(W)$ to be the exterior algebra. Set
\begin{equation*}
\det W =\bw{\dim W} (W).
\end{equation*}
Clearly, $\det W$ is a complex line. If $W=\{0\}$, then
\begin{equation*}
\det
W=\bC.
\end{equation*}
Let $E^{\sbullet}=\bigoplus_{i\in \bZ} E^{i}$ be a finite dimensional $\bZ$-graded
space. Put
\begin{equation*}
\det E^{\sbullet}= \bigotimes_{i\in \bZ} (\det E^{i})^{(-1)^{i}}.
\end{equation*}

For $m\in \bN$, let
\begin{equation*}
(C^{\sbullet},d): 0\to C^{0}\to C^{1}\to\cdots\to C^{m}\to 0
\end{equation*}
be a complex of finite dimensional vector
spaces. By \cite{KM76} or
\cite[(1.5)]{BGS1}, we
have the canonical isomorphism of lines
\begin{equation}\label{detCdetH}
\tau_{C^{\sbullet}}: \det C^{\sbullet}\simeq \det H^{\sbullet}(C^{\sbullet},d).
\end{equation}
If $s_{j}^{k}\in C^{k}$ such that $\{s_{j}^{k}\}_{j=1}^{k_{j}}$ projects to a basis of
$C^{k}/\ker (d_{|C^{k}})$, if $\mu_{j}^{k}\in \ker( d_{|C^{k}})$ such that
$\{\mu_{j}^{k}\}_{j=1}^{k_{j}'}$ projects to a basis of
$H^{k}(C^\sbullet,d)$,
then
\begin{equation}\label{Maja1}
(\wedge_{j} s^{0}_{j}\otimes \wedge_{j}\mu^{0}_{j})\otimes
(\wedge_{j} (ds^{0}_{j})\otimes
\wedge_{j}s^{1}_{j} \otimes \wedge_{j}\mu^{1}_{j})^{-1} \otimes\cdots \otimes(\wedge_{j}
(ds^{m-1}_{j})\otimes \wedge_{j}\mu^{m}_{j})^{(-1)^m}
\end{equation}
defines a canonical element of $\det C^\sbullet$.
If $\ol{\mu}_{j}^{k}$ denotes the image of $\mu_{j}^{k}$ in $H^{k}(C^{\sbullet},d)$, under
\eqref{detCdetH}, the element \eqref{Maja1} maps to
\begin{equation}\label{Maja11}
( \wedge_{j}\ol{\mu}^{0}_{j})\otimes
( \wedge_{j}\ol{\mu}^{1}_{j})^{-1} \otimes\cdots \otimes(
\wedge_{j}\ol{\mu}^{m}_{j})^{(-1)^m}\in \det
H^{\sbullet}(C^{\sbullet},d).
\end{equation}

\subsection{The cohomology of $\bbS^{1}$}\label{sCS1}
Let $\bbS^{1}=\bR/\bZ$ be an oriented circle. Let $F$ be a flat vector
bundle of rank $r$ on $\bbS^{1}$. Let $\rho: \pi_{1}(\bbS^{1})\to
\GL_{r}(\bC)$ be the holonomy\footnote{For any flat vector bundle
$F$ on a manifold $X$, the holonomy is a representation $\rho:
\pi_{1}(X)\to \GL_{r}(\bC)$ of the fundamental group
$\pi_{1}(X)$ such that $F=\pi_{1}(X)\setminus (\widetilde{X}\times
\bC^{r})$, where $\widetilde{X}$ is the universal cover of $X$ and $\pi_{1}(X)$ acts on the left on $\widetilde{X}$ by
the deck transformation
and on $\bC^{r}$ by $\rho$.} of $F$. Let $a_{0}\in \pi_{1}(\bbS^{1})$ be the
generator of $\pi_{1}(\bbS^{1})$, which is compatible with the
orientation on~$\bbS^{1}$. Set $A=\rho(a_{0})\in
\GL_{r}(\bC)$.

Consider the canonical triangulation on $\bbS^{1}$ induced by one $0$-simplex
$\sigma_{0}$ and one $1$-simplex $\sigma_{1}$ as in Figure \ref{fig:0}.
\begin{figure}[htbp]
\centering
\includegraphics[width=1in]{shen-yu_fig1}
\caption{A triangulation on $\bbS^{1}$.}\label{fig:0}
\end{figure}
It induces a complex of simplicial cochains with values in $F$ given
by
\begin{equation*}\label{CS1}
(C^{\sbullet}(\bbS^{1},F),d):0 \to\bC^{r}\To{A-1}\bC^{r} \to0.
\end{equation*}
By \eqref{detCdetH}, the canonical element $
\mathbf{1}\in \bC=\det C^{\sbullet}(\bbS^{1},F)$ defines an element
\begin{equation}\label{eqsigmaa}
\sigma_{A}\in \det H^{\sbullet}(\bbS^{1},F).
\end{equation}
We equip $\det H^{\sbullet}(\bbS^{1},F)$ with a metric
$\|{\cdot}\|^{2}_{\det H^{\sbullet}(\bbS^{1},F)}$ such that
\begin{equation}\label{nsa}
\|\sigma_{A}\|_{\det H^{\sbullet}(\bbS^{1},F)}=1.
\end{equation}

If $1$ is not an eigenvalue of $A$, then $H^{\sbullet}(\bbS^{1},F)=0$.
By \eqref{Maja1}, the norm of the canonical element
$\mathbf{1}\in \bC=\det H^{\sbullet}(\bbS^{1},F)$ is given by
\begin{equation}\label{sa}
\|\mathbf{1}\|_{\det
H^{\sbullet}(\bbS^{1},F)}=\left|{\det\,}(1-A)\right|^{-1}.
\end{equation}

\begin{re}
Equation \eqref{sa} is just Proposition \ref{prop} for the rotation flow on $\bbS^{1}$.
\end{re}

\begin{re}
Since the flat vector bundle is not necessarily unimodular, \ie $|{\det\,}(A)|$ is
not necessarily equal to $1$, the choice of the
orientation on $\bbS^{1}$ is very important.
\end{re}

\subsection{A fusion principle}\label{Fp}
Let $Y_{1}\subset Y_{2}\subset Y_{3}$ be three compact smooth
manifolds with boundaries of
the same dimension such that $Y_{1}\subset \mathring{Y_{2}}$ and
$Y_{2}\subset \mathring{Y_{3}}$. Let $F$ be a flat vector bundle on
$Y_{3}$ of rank $r$. Denote by $H^{\sbullet}(Y_{3},Y_{2},F)$,
$H^{\sbullet}(Y_{3},F), \ldots,$ the corresponding relative or absolute
cohomologies with coefficients in $F$.

As in \cite[(0.16)]{BruningMa_2012}, we have a long exact sequence
\begin{multline}\label{Les}
\cdots\to H^{i}(Y_{3},Y_{2},F)\to H^{i}(Y_{3},Y_{1},F)\to
H^{i}(Y_{2},Y_{1},F)\\
\to H^{i+1}(Y_{3},Y_{2},F) \to \cdots.
\end{multline}
By \eqref{detCdetH} and \eqref{Les}, we get an isomorphism of lines
\begin{equation}\label{Maja3}
f_{21,32 }:\det H^{\sbullet}(Y_{2},Y_{1},F)\otimes \det H^{\sbullet}(Y_{3},Y_{2},F)\simeq \det H^{\sbullet}(Y_{3},Y_{1},F).
\end{equation}
Using the other triples $(\varnothing,Y_{1},Y_{2})$ and
$(\varnothing,Y_{1},Y_{3})$, we get similar isomorphisms
\begin{equation}\label{Maja4}
\begin{aligned}
f_{1,21}: \det H^{\sbullet}(Y_{1},F)\otimes \det
H^{\sbullet}(Y_{2},Y_{1},F)\simeq \det H^{\sbullet}(Y_{2},F),\\
f_{1,31}: \det H^{\sbullet}(Y_{1},F)\otimes \det
H^{\sbullet}(Y_{3},Y_{1},F)\simeq \det H^{\sbullet}(Y_{3},F).
\end{aligned}
\end{equation}
By \eqref{Maja3} and \eqref{Maja4}, we see that
$f_{2,32}\circ(f_{1,21}\otimes\id) $ and $f_{1,31}\circ (\id \otimes f_{21,32})$
define two isomorphisms
\begin{equation}\label{FYYY}
\det H^{\sbullet}(Y_{1},F)\otimes \det H^{\sbullet}(Y_{2},Y_{1},F)\otimes \det
H^{\sbullet}(Y_{3},Y_{2},F) \simeq \det H^{\sbullet}(Y_{3},F).
\end{equation}

\begin{prop}\label{propFp}
There is $\mu=1$ or $-1$\footnote{See \cite[Rem.\,1.2]{BGS1} or
\cite{KM76} for the detail about the sign. } such that
\begin{equation}\label{mabu}
f_{2,32} \circ(f_{1,21}\otimes\id) = \mu \times
f_{1,31}\circ (\id \otimes f_{21,32}).
\end{equation}
\end{prop}
\begin{proof}
As in \cite[(0.15)]{BruningMa_2012}, let us take a smooth triangulation
of $Y_{3}$ such that it induces also smooth triangulations on
$Y_{1}$ and $Y_{2}$. Denote by
\[
(C^{\sbullet}(Y_{1},F),d),\;
(C^{\sbullet}(Y_{2},Y_{1},F),d), \ldots,
\]
the complexes of simplicial
cochains with coefficients in $F$. Then we have an exact sequence of
complexes
\begin{equation}\label{eq111}
0\to (C^{\sbullet}(Y_{2},Y_{1},F),d)\to (C^{\sbullet}(Y_{2},F),d)\to
(C^{\sbullet}(Y_{1},F),d)\to 0.
\end{equation}
By \eqref{detCdetH} and \eqref{eq111}, we get an isomorphism of lines
\begin{equation*}
f^{C}_{1,21}: \det C^{\sbullet}(Y_{1},F)\otimes \det
C^{\sbullet}(Y_{2},Y_{1},F)\to \det C^{\sbullet}(Y_{2},F).
\end{equation*}
We can define $f^{C}_{2,32}$, $f^{C}_{1,31}$ and $f^{C}_{21,32}$ in a
similar way. By an easy calculation, there is $\mu=1$ or $-1$ such that
\begin{equation}\label{mabuC}
f_{2,32}^{C} \circ(f_{1,21}^{C}\otimes\id) = \mu \times
f_{1,31}^{C}\circ (\id \otimes f_{21,32}^{C}).
\end{equation}

By \eqref{Maja1} and \eqref{Maja11}, there is $\mu=1$ or $-1$ such that the diagram
\begin{equation}\label{eqBGSMA}
\begin{aligned}
\xymatrix{
\det C^{\sbullet}(Y_{1},F)\otimes \det C^{\sbullet}(Y_{2},Y_{1},F)
\ar[d]^{\tau_{C^{\sbullet}(Y_{1},F)}\otimes\tau_{C^{\sbullet}(Y_{2},Y_{1},F)}}
\ar[r]^-{f^{C}_{1,21}} &\det C^{\sbullet}(Y_{2},F)
\ar[d]^{\mu\tau_{C^{\sbullet}(Y_{2},F)}}
\\
\det H^{\sbullet}(Y_{1},F)\otimes \det
H^{\sbullet}(Y_{2},Y_{1},F)\ar[r]^-{f_{1,21}} &\det H^{\sbullet}(Y_{2},F)}
\end{aligned}
\end{equation}
commutes. Tensoring each vertical line of \eqref{eqBGSMA} by the
isomorphism
\begin{equation*}
\tau_{C^{\sbullet}(Y_{3},Y_{2},F)}:\det C^{\sbullet}(Y_{3},Y_{2},F)\simeq
\det H^{\sbullet}(Y_{3},Y_{2},F),
\end{equation*}
and using \eqref{eqBGSMA} again for the pair $(Y_{2},Y_{3})$, we
find that
there is $\mu=1$ or $-1$ such that the diagram
\begin{equation}\label{eqDjht}
\hbox{\smaller$
\let\labelstyle\scriptstyle
\xymatrixcolsep{6pc}\xymatrix{
\det C^{\sbullet}(Y_{1},F)\otimes \det
C^{\sbullet}(Y_{2},Y_{1},F)\otimes \det C^{\sbullet}(Y_{3},Y_{2},F)
\ar[d]^{\tau_{C^{\sbullet}(Y_{1},F)}\otimes\tau_{C^{\sbullet}(Y_{2},Y_{1},F)}\otimes\tau_{C^{\sbullet}(Y_{3},Y_{2},F)}}
\ar[r]^-{f^{C}_{2,32} \circ(f^{C}_{1,21}\otimes\id)
} &\det C^{\sbullet}(Y_{3},F)
\ar[d]^{\mu\tau_{C^{\sbullet}(Y_{3},F)}}
\\
\det H^{\sbullet}(Y_{1},F)\otimes \det
H^{\sbullet}(Y_{2},Y_{1},F)\otimes \det
H^{\sbullet}(Y_{3},Y_{2},F)\ar[r]^-{f_{2,32} \circ(f_{1,21}\otimes\id) }
&\det H^{\sbullet}(Y_{3},F)}$}
\end{equation}
commutes.
In \eqref{eqDjht}, if we replace the
horizontal morphisms by $f_{1,31}^{C}\circ (\id \otimes
f_{21,32}^{C})$ and $f_{1,31}\circ (\id \otimes
f_{21,32})$, the corresponding diagram still commutes. By \eqref{mabuC}, we get~\eqref{mabu}.
\end{proof}

\section{Milnor metric}\label{S1M}
This section is organized as follows. In Sections \ref{sMS} and
\ref{sDZf}, we recall
the definitions of Morse-Smale flow and the associated Ruelle dynamical zeta
function. In Section~\ref{sMSMF}, we recall some results due to Franks
\cite{Franks1979,Franks1982} on the construction of a new
gradient flow by destroying the closed orbits of the original Morse-Smale flow. In~Section~\ref{sSF}, using the Smale filtration, we introduce the Milnor
metric. In~Section~\ref{sCMM}, we establish a comparison formula for the two Milnor
metrics, one is associated with the Morse-Smale flow and the other is
associated with the gradient flow constructed by Franks.

We refer the reader to the classical textbook of Palis and de
Melo \cite{Palis82} for the basic notion on dynamical system.

\subsection{Morse-Smale flow}\label{sMS}
Let $X$ be a connected closed smooth manifold of dimension $m$. Let $V$ be a
vector field on $X$. Consider the differential equation
\begin{equation}\label{eqDiffV}
\frac{dx}{dt}=V(x).
\end{equation}
Equation \eqref{eqDiffV} defines a group of diffeomorphism
$(\phi_{t})_{t\in \bR}$ of $X$.

If $x\in X$, an orbit of $x$ is defined by the image $t\in \bR\mto
\phi_{t}(x)\in X$. We call $x\in X$ is a fixed point, if its orbit reduces to a point, i.e, for all $t\in \bR$,
\begin{equation*}
\phi_{t}(x)=x.
\end{equation*}
Clearly, $x\in X$ is a fixed point if and only if $V(x)=0$.
We call an orbit is closed if it is diffeomorphic to $\bbS^{1}$.
Denote by $A$ the set of fixed points and by $B$ the set of closed orbits.

\begin{defin}
A fixed point $x\in X$ of the flow $\phi_{\cdot}$ is called
hyperbolic if there is a $\phi_{t}$-invariant splitting\vspace*{-3pt}
\begin{equation*}
T_{x}X= E^{\ru}_{x}\oplus E^{\rs}_{x},
\end{equation*}
and there exist $C>0,\theta>0$ and a Riemannian metric $g^{TX}$ on $X$
such that for $v\in E^{\ru}_{x}$, $v'\in E^{\rs}_{x}$, and $t>0$, we have\vspace*{-3pt}
\begin{equation*}\label{eqMS}
\left|\phi_{- t,*}v\right|\l
Ce^{-\theta t}\left|v\right|,\quad \left|\phi_{ t,*}v'\right|\l
Ce^{-\theta t}\left|v'\right|.
\end{equation*}
The unstable and stable manifolds of the hyperbolic fixed point $x$ are defined by
\begin{equation*}
W^{\ru}_{x}=\Big\{y\in X: \lim_{t\to -
\infty}d_{X}(\phi_{t}(y),x)=0\Big\},\quad W^{\rs}_{x}= \Big\{y\in X: \lim_{t\to
\infty}d_X(\phi_{t}(y),x)=0 \Big\},
\end{equation*}
where $d_{X}$ denotes the Riemannian distance on $(X,g^{TX})$. The index $\ind(x)\in \bN$ of $x$ is defined by\vspace*{-3pt}
\begin{equation}\label{eqindxEu}
\ind(x)=\dim E^{\ru}_{x}.
\end{equation}
\end{defin}

Note that if $V=-\nabla f$ is a negative gradient of a Morse function
$f$
with respect to some Riemannian metric, then the index $\ind(x)$ of
the critical point $x$ is just the Morse index of $f$ at $x$.

\begin{defin}
A closed orbit $\gamma$ of the flow $\phi_{\cdot}$ is called hyperbolic,
if there is a $\phi_{t}$-invariant continuous splitting\vspace*{-3pt}
\begin{equation*}\label{eq:AF1}
TX|_{\gamma}=\bR V\oplus E^{\ru}_{\gamma}\oplus E^{\rs}_{\gamma},
\end{equation*}
of $C^{0}$-vector
bundles over $\gamma$ such that \eqref{eqMS} holds. The associated unstable and
stable manifolds are defined by
\begin{equation*}
\begin{aligned}
W^{\ru}_{\gamma}&=\bigcup_{x\in \gamma} \bigl\{y\in X: \textstyle\lim_{t\to
-\infty}d_{X}(\phi_{t}(y),\phi_{t}(x))=0 \bigr\},\\
W^{\rs}_{\gamma}&=\bigcup_{x\in \gamma} \bigl\{y\in X: \textstyle\lim_{t\to
+\infty}d_{X}(\phi_{t}(y),\phi_{t}(x))=0 \bigr\}.
\end{aligned}
\end{equation*}
The index $\ind(\gamma)\in \bN$ of $\gamma$ is defined by
\begin{equation}\label{indexg}
\ind(\gamma)=\rk E^{\ru}_{\gamma}.
\end{equation}
\end{defin}

\begin{defin}
The nonwandering set of $\phi_{\cdot}$ is defined by\vspace*{-3pt}
\begin{equation*}\textstyle
\bigl\{x\in X: \hbox{$\forall$ open neighborhood $U$ of $x$,
$\forall\, T>0$, } U\cap\bigcup_{t\g T} \phi_{t}(U)\neq
\varnothing \bigr\}.
\end{equation*}
\end{defin}
Clearly, $A\cup \bigcup_{\gamma\in B}\gamma$ is contained in the nonwandering set.

\begin{defin}
A vector field $V$ or a flow $\phi_{\cdot}$ is called Morse-Smale if
\begin{itemize}
\item the sets $A$ and $B$ are finite and contain only hyperbolic elements;
\item the nonwandering set of
$\phi_{\cdot}$ is equal to
$ A\cup \bigcup_{\gamma\in B}\gamma$;
\item the stable and unstable manifold of any critical
element in $A\coprod B$ intersect transversally.
\end{itemize}
\end{defin}

In the sequel, we assume
that $V$ is a
Morse-Smale vector field.

\subsection{Ruelle dynamical zeta function}\label{sDZf}
For $\gamma\in B$, denote by $\ell_{\gamma}\in \bR^{*}_{+}$ its
minimal period. A closed orbit $\gamma\in B$
is called untwist if $E^{\ru}_{\gamma}$ is orientable along $\gamma$, and is
called twist otherwise. Put
\begin{equation}\label{Delta}
\Delta(\gamma)=
\begin{cases}
1&\hbox{if } \gamma \hbox{ is untwist},\\
-1&\hbox{if } \gamma \hbox{ is twist}.
\end{cases}
\end{equation}

Let $\rho:\pi_{1}(X)\to \GL_{r}(\bC)$ be a representation of the fundamental group of $X$. If~$\gamma\in B$, denote by $\rho(\gamma)$
the holonomy along $\gamma$. Clearly, $\rho(\gamma)$ is well-defined up to a conjugation.

\begin{defin} The twist Ruelle dynamical zeta function is a meromorphic
function on $\bC$ defined for $s\in \bC$ by
\begin{equation}\label{rrs}
R_{\rho}(s)=\prod_{\gamma\in
B}\det(1-\Delta(\gamma)\rho(\gamma)e^{-s
\ell_\gamma})^{(-1)^{\ind(\gamma)}}.
\end{equation}
\end{defin}

\subsection{Franks' Morse function}\label{sMSMF}
We follow \cite[\S 1]{Franks1979}. Let $\bbD^{r}$ be the
$r$-dimensional open unit ball of center $0\in
\bR^{r}$.
A fixed point $x\in A$ of index $p$ is said to be of standard form if there is a system of coordinates $(y_{1},\ldots, y_{m})\in \bbD^{m}$ on a neighborhood
of $x$ such that $x$ is represented by $0$ and
\begin{equation*}
V=\sum_{i=1}^{p}y_{i}\frac{\p}{\p
y_{i}}-\sum_{i=p+1}^{m}y_{i}\frac{\p}{\p y_{i}}.
\end{equation*}

For closed orbits we must distinguish the following four cases in
establishing the standard forms. Assume $\gamma\in B$ such that
$\ind(\gamma)=p$.

\Case\label{Case:1}
Suppose that $TX|_{\gamma}$ is orientable and that $\gamma$ is
untwist. In this case, $\gamma$ is said to be of standard form, if there
is a system of coordinates $(t,y_{1},\ldots,
y_{m-1})\in \bbS^{1}\times \bbD^{m-1}$
on a tubular neighborhood
$U_{\gamma}$ of $\gamma$ such that $\gamma$ is represented by
$(t,0)\in \bbS^{1}\times \bbD^{m-1}$ and
\begin{equation*}\label{eqV}
V=\frac{\p}{\p t}+\sum_{i=1}^{p}y_{i}\frac{\p}{\p
y_{i}}-\sum_{i=p+1}^{m-1}y_{i}\frac{\p}{\p y_{i}}.
\end{equation*}

\Case\label{Case:2}
Suppose that $TX|_{{\gamma}}$ is orientable and
that $\gamma$ is
twist. In this case, $\gamma$ is said to be of standard form, if
$U_{\gamma}$ and $V$ can be obtained
from Case \ref{Case:1} by the identification
\begin{equation*}\label{enid}
(t,
x_{1},\ldots,x_{m-1})\sim (t+1/2,
-x_{1},x_{2},\ldots, x_{p}, -x_{p+1}, x_{p+2},\ldots,x_{m-1}).
\end{equation*}

\Case\label{Case:3}
Suppose that $TX|_{{\gamma}}$ is not orientable and
that $\gamma$ is
untwist. In this case, $\gamma$ is said to be of standard form, if
$U_{\gamma}$ and $V$ can be obtained
from Case \ref{Case:1} by the identification
\begin{equation*}
(t,
x_{1},\ldots,x_{m-1})\sim (t+1/2,
x_{1},\cdots,x_{p},-x_{p+1},x_{p+1},\ldots,x_{m-1}). \end{equation*}

\Case\label{Case:4}
Suppose that $TX|_{{\gamma}}$ is not orientable
and that $\gamma$ is
twist. In this case, $\gamma$~is said to be of standard form, if
$U_{\gamma}$ and $V$ can be obtained
from Case \ref{Case:1} by the identification
\begin{equation*}
(t,
x_{1},\ldots,x_{m-1})\sim (t+1/2,
-x_{1},x_{2},\cdots,x_{m-1}).
\end{equation*}
Note that in \cite[\S 1]{Franks1979} the author assumed that $X$ is orientable, so only the first two cases appear.

The following three propositions are \cite[Prop.\,1.6,
Th.\,2.2, Prop.\,5.1]{Franks1979}.
Their proofs can be generalized to the non orientable case
with some evident modifications. We omit the details.

\begin{prop}\label{propV01}
For any Morse-Smale vector field $V$, there is a smooth family of Morse-Smale
vector fields $(V_{\ell})_{0\l \ell\l 1}$ such that $V_{0}=V$ and that the critical
elements of~$V_{1}$ are all of standard forms and are precisely the same as the critical
elements of~$V$. Moreover, $V_0$ and $V_1$ are topologically
conjugated, \ie there is a homeomorphism carrying the orbits
of $V_0$ to those of $V_{1}$ and preserving their orientations.
\end{prop}

\begin{re}
The second part of Proposition \ref{propV01} is a consequence of
the Structural stability of the Morse-Smale flow \cite{Palis68,PalisSmale70}.
\end{re}

\begin{re}\label{reVl}
Following the proof of \cite[Prop.\,1.6]{Franks1979}
given by Franks, we can choose the family $(V_{\ell})_{0\l
\ell \l1}$ such that the critical elements are preserved
under the deformation. However, in the proof of our main result
Theorem \ref{thm1} given in Section~\ref{S2}, we need only
choose a family such that all the set of the fixed points of
$V_{\ell}$ are in a small neighbourhood of the set of the fixed points of $V$.
\end{re}

The relation between the Morse-Smale flow and the Morse function is
summarized in the following two propositions. The first one is due to
Smale \cite[Th.\,B]{S61}.

\begin{prop}
If $V$ is a Morse-Smale vector field whose flow has fixed points
in standard form and no closed orbits, then $V$ is a negative gradient of
a certain Morse function with respect to some Riemannian metric.
\end{prop}

To state the following proposition, let us introduce some notation.
For $x,y\in A$ such that $\ind(y)=\ind(x)-1$, then $W^{\ru}_{x}\cap W^{\rs}_{y}$
consists of a finite set $\Gamma(x,y)=\{a_{\sbullet}\}$ of integral curves of
$V$ such that $a_{-\infty}=x$ and $a_{\infty}=y$ along
which $W^{\ru}_{x}$ and $W^{\rs}_{y}$ intersect transversally. Let us fix
an orientation on each $W^{\ru}_{x}$ with $x\in A$. Define
$n(a)=\pm1$ as in \cite[(1.28)]{BZ92}, whose definition does not
require the manifold to be orientable.

\begin{prop}\label{propfV}
For some small neighborhood $U=\bigcup_{\gamma\in B}U_{\gamma}$ of
closed orbits $\bigcup_{\gamma\in B}\gamma$, there is
a Morse function $f$ on $X$ whose gradient vector field $\nabla f$
with respect to a certain Riemannian metric is
Morse-Smale, such that
\begin{itemize}
\item on $X\setminus U$, we have
\begin{equation}\label{eqdf=V}
-\nabla f=V,
\end{equation}
\item on each $U_{\gamma}$, the Morse function $f$ has only two critical points $x_{\gamma},x'_{\gamma}$ of
index $\ind(\gamma)+1$ and $\ind(\gamma)$ respectively.
\end{itemize}
Also, $\Gamma(x_{\gamma},x'_{\gamma})$ consists of two
integral curves $a_{\gamma},a'_{\gamma}$ (see Figure \ref{fig:1})
such that their composition $a_{\gamma}\circ (a'_{\gamma})^{-1}$ and the
closed orbit $\gamma$ lie in the same freely homotopy class of
loops on $X$ and that\footnote{This
requires a choice of the orientations on the unstable manifolds of
$x_{\gamma},x_{\gamma}'$. Such choice is irrelevant.}
\begin{equation*}\label{nana}
n(a_{\gamma})n(a'_{\gamma})=-\Delta(\gamma).
\end{equation*}
\end{prop}

\begin{figure}[htbp]
\vspace*{-.5\baselineskip}
\centering
\includegraphics[width=3.5in]{shen-yu_fig2}
\vspace*{-.5\baselineskip}
\caption{A closed orbit and integral curves }\label{fig:1}
\end{figure}

\begin{re}\label{reK}We recall the essential step in the
construction of the gradient $\nabla f$ given by Franks \cite[Prop.\,8.5]{Franks1982}. For simplicity, assume that there is only one
closed orbit $\gamma$ and it is of index $p$ and in standard
form of Case \ref{Case:1}.
For $\delta>0$ small enough, let $\rho\in
C^{\infty}_{c}(\bbD^{m-1},[0,1])$
be a cutoff function, which is equal to
$1$ on $|y|<\delta$ and to $0$ on $|y|>2\delta$. Put
\begin{equation*}
V_{1}=
\begin{cases}
\displaystyle
\big((1-\rho)+\rho\sin(2\pi t)\big)\frac{\p}{\p t}+\sum_{i=1}^{p}y_{i}\frac{\p}{\p
y_{i}}-\sum_{i=p+1}^{m-1}y_{i}\frac{\p}{\p y_{i}}, &\hbox{ in
} U_{\gamma},\\
V,& \hbox{ outside } U_{\gamma}.
\end{cases}
\end{equation*}
It is easy to see that the
nonwandering set of $V_{1}$ in $U_{\gamma}$ consists of two
points $(1,0),(\sfrac{1}{2},0)\in
\bbS^{1}\times \bbD^{m-1}$. Then, by a small perturbation on $V_{1}$, we get a Morse-Smale gradient vector
field $-\nabla f$ which
has the desired transversality and other properties.

We remark that by the above construction, we can find a family of vector fields $(V_{\e})_{0\l\e\l1}$ connecting $V$ and
$-\nabla f$ such that near $\{1/4\}\times \bbD^{m-1}$, for
any $ \e\in [0,1]$, $V_{\e}$~does not vanish.
Similar remark holds for $\gamma$ in standard forms of Cases \ref{Case:2}--\ref{Case:4}.
In~Section~\ref{s33}, we will use this fact to simplify the proof of our main theorem.
\end{re}

\subsection{Smale filtration and Milnor metric}\label{sSF}
Following \cite[Def.\,9.10]{Franks1982},
let
\begin{equation}\label{eqSmale}
\varnothing=X^{0}\subset X^{1}\subset\cdots \subset X^{N}=X
\end{equation}
be a Smale filtration on $X$ associated with $V$. Note that each
$X^{p}\subset X$ is a submanifold with boundary, and can be constructed by the sublevel set of a smooth Lyapunov function. Also, we
have
\begin{itemize}
\item on each $\p X^{p}$, $V$ does not vanish and points toward the inside of $X^{p}$;
\item there is only one critical element $c\in A\coprod
B$ contained in $X^{p+1}\setminus
X^{p}$ and
\begin{equation*}
\{c\}=\bigcap_{t\in \bR}\phi_{t}(X^{p+1}\setminus
X^{p}).
\end{equation*}
\end{itemize}

Let $(F,\nabla^{F})$ be a flat vector bundle on $X$ induced by the representation
$\rho$.
Let $H^{\sbullet}(X,F)$ be the cohomology of the sheaf of locally constant
sections of $F$. Put
\begin{equation}\label{eqlambda}
\lambda=\det H^{\sbullet}(X,F).
\end{equation}
We use the notation in Section \ref{Fp}. By \eqref{FYYY}, we get an isomorphism
\begin{equation}\label{eqSS}
\sigma_{V}:\bigotimes_{p=0}^{N-1}\det
H^{\sbullet}(X^{p+1},X^{p},F)\simeq \lambda.
\end{equation}
By Proposition \ref{propFp}, up to a sign, the
morphism $\sigma_{V}$ does not depend on the way that the
cohomologies are fused.

By \cite[Th.\,9.11]{Franks1982} (see also \cite[\S 2]{MorgadoMorseSmale}), if the critical element $c\in
X^{p+1}\setminus X^{p}$ is a fixed
point, then
\begin{equation}\label{E1x}
H^{q}(X^{p+1},X^{p},F)=
\begin{cases}
F_{c},& q=\ind (c),\\
0,&\hbox{otherwise},
\end{cases}
\end{equation}
and if the critical element $c\in
X^{p+1}\setminus X^{p}$ is a closed orbit, then
\begin{equation}\label{E1g}
H^{q}(X^{p+1},X^{p},F)=
\begin{cases}
H^{q-\ind (c)}\big(c,o(E^{\ru}_{c})\otimes F|_{c}\big),& q-\ind( c)=0
\hbox{ or }1,\\
0,&\hbox{otherwise,}
\end{cases}
\end{equation}
where $o(E_{c}^{\ru})$ is the orientation line bundle of $E^{\ru}_{c}$ along the
closed orbit $c$.

We equip $\det
H^{\sbullet}(\gamma,o(E^{\ru}_{\gamma})\otimes
F|_{\gamma})$ with the metric $\|{\cdot}\|^{2}_{\det
H^{\sbullet}(\gamma,o(E^{\ru}_{\gamma})\otimes
F|_{\gamma})}$ defined in~\eqref{nsa}.
Let $g^{F}$ be a Hermitian metric on $F$. By \eqref{eqSS}--\eqref{E1g}, the
restriction $g^{F}|_{A}$
and the metric $\|{\cdot}\|^{2}_{\det
H^{\sbullet}(\gamma,o(E^{\ru}_{\gamma})\otimes
F|_{\gamma})}$ induce a metric
$\|{\cdot}\|^{\rM,2}_{\lambda,V}$ on
$\lambda$. By Proposition \ref{propFp},
this metric does not depend on the way that the
cohomologies are fused.

\begin{defin}
The metric $\|{\cdot}\|^{\rM,2}_{\lambda,V}$ on
$\lambda$ is called the Milnor metric associated with $V$.
\end{defin}

\begin{re}
If $V=-\nabla f$ is a negative gradient of a Morse function $f$, then
$\|{\cdot}\|^{\rM,2}_{\lambda,V}$ coincides with the one constructed by
Bismut-Zhang \cite[Def.\,1.9]{BZ92}.
In fact, there is a small difference with Bismut-Zhang's
construction, where they used a filtration \cite[(1.37)]{BZ92} induced
by sublevel sets of a nice Morse function. Using Proposition~\ref{propFp}, we can deduce that the two constructions coincide.
\end{re}

\begin{re}\label{reDMA}
For two topologically conjugated Morse-Smale vector fields whose
critical elements coincide, we can choose the same Smale filtration. From our construction, the corresponding Milnor metrics coincide.
\end{re}

\begin{re}
The Milnor metric for
general Morse-Smale flow does not depend on the Smale
filtration \eqref{eqSmale}. We will not give a direct proof since it is a consequence of our main Theorem \ref{thm1}.
\end{re}

Let us restate and prove Proposition \ref{prop}.

\begin{prop}\label{Prop}
If $V$ does not have any fixed points, and
if none of $\Delta(\gamma)$ is an eigenvalue of $\rho(\gamma)$, then
$H^{\sbullet}(X,F)=0$, and the norm of the canonical section
$\mathbf{1}\in \bC=\det H^{\sbullet}(X,F)$ is given by
\begin{equation}\label{eqReMi}
\|\mathbf{1}\|^{\rM}_{\lambda,V}=|R_{\rho}(0)|^{-1}.
\end{equation}
\end{prop}
\begin{proof}
For $\gamma\in B$, the holonomy of $o(E_{\gamma}^{\ru})\otimes
F|_{\gamma}$
along $\gamma$ is $\Delta(\gamma)\rho(\gamma)$. By our
assumption,
\begin{equation}\label{eqDMA}
H^{\sbullet}(\gamma,o(E_{\gamma}^{\ru})\otimes
F_{\gamma})=0.
\end{equation}
By \eqref{Les}, \eqref{E1x}, \eqref{E1g}, and \eqref{eqDMA}, we can deduce that $H^{\sbullet}(X,F)=0$. By
\eqref{sa}, \eqref{rrs}, and \eqref{eqSS}, we get
\eqref{eqReMi}.
\end{proof}

\subsection{A comparison formula for Milnor metrics}\label{sCMM}
In this section, we assume that all the critical elements of
$V$ are in standard forms, and that $f$ is chosen as in
Proposition \ref{propfV}.

Let $\det \tau({a'_{\gamma}})\in \det
F_{x'_{\gamma}}\otimes (\det
F_{x_{\gamma}})^{-1}$ be the canonical element induced by the
parallel transport with respect to the flat connection along the
integral curve $a_{\gamma}'$ (see Figure \ref{fig:1}). Let $\|{\cdot}\|^{2}_{\det
F_{x'_{\gamma}}\otimes (\det
F_{x_{\gamma}})^{-1}}$ be the metric on $\det
F_{x'_{\gamma}}\otimes (\det
F_{x_{\gamma}})^{-1}$ induced by $g^{F}_{x_{\gamma}}$ and
$g^{F}_{x'_{\gamma}}$. 

\begin{prop}\label{propMM}
The following identity holds,
\begin{equation}\label{eqMM}
\log\bigl(\sfrac{\|{\cdot}\|^{\rM,2}_{\lambda,-\nabla
f}}{\|{\cdot}\|^{\rM,2}_{\lambda,V}}\bigr)=\sum_{\gamma\in
B}(-1)^{\ind(\gamma)}\log
\left\|{\det\,}\tau(a_{\gamma}')\right\|^{2}_{\det
F_{x'_{\gamma}}\otimes (\det
F_{x_{\gamma}})^{-1}}.
\end{equation}
\end{prop}
\begin{proof}
We refine the filtration \eqref{eqSmale} by the new critical
points of $f$. By Propositions \ref{propFp} and \ref{propfV}, and by
\eqref{eqSS}, the following diagram commutes
\begin{equation}\label{eqDia}
\begin{aligned}
\xymatrix{
\bigotimes_{x\in A} \big(\!\det F_{x}\big)^{ (-1)^{\ind (x)}}
\bigotimes_{\gamma\in B} \left\{\det
F_{x'_{\gamma}}\otimes (\det
F_{x_{\gamma}})^{-1}\right\}^{(-1)^{\ind (\gamma)}}
\ar@<-1cm>[d]
\ar[r]^-{\sigma_{-\nabla f}} &\lambda\ar[d]
\\
\bigotimes_{x\in A}\!\big(\!\det
F_{x}\big)^{(-1)^{\ind (x)}} \!
\bigotimes_{\gamma \in B}\Big\{\!\det
H^{\sbullet}(\gamma,o(E^{\ru}_{\gamma})\otimes
F|_{\gamma})\!\Big\}^{(-1)^{\ind
(\gamma)}}\ar[r]^-{\sigma_V} &\lambda}
\end{aligned}
\end{equation}
where the first vertical arrow is induced by \eqref{eqsigmaa} and
the second vertical arrow is a multiplication by $\pm 1$. The Milnor metric $\|{\cdot}\|_{\lambda,-\nabla f}^{\rM,2}$ is
obtained from the metric $g^{F}|_{A \cup
\{x_{\gamma},x'_{\gamma}:\gamma\in B\}}$ via $\sigma_{-\nabla f}$.
By \eqref{eqDia}, we get
\eqref{eqMM}.
\end{proof}

\section{An extension of Bismut-Zhang's theorem to Morse-Smale
flow}\label{S2}
This section is organized as follows. In Sections
\ref{sec:Bere}-\ref{sMa}, following \cite{BZ92}, we recall the
constructions of the Berezin integral, the Mathai-Quillen current, and the Ray-Singer metric.
In Section \ref{s33}, we restate and prove our main theorem.
\subsection{Berezin integral}\label{sec:Bere}
Let $E$ be a real Euclidean space of dimension $n$ with
the scalar product $\<,\>$, and let $W$ be a real vector space of
finite dimension. We use the supersymmetric formalism of Quillen
\cite{Quillensuper}. Denote by $\wotimes$ the tensor product of super algebras.

Suppose temporarily that $E$ is oriented and that $e_1,\ldots, e_n$
is an oriented orthonormal basis of $E$. Let $e^1,\ldots,e^n$ be the corresponding dual basis of $E^*$. We~define $\int^B$ to be the linear map from $\bw\sbullet(W^*)\wotimes\bw\sbullet(E^*)$ into $\bw\sbullet(W^*)$, such that if $\alpha\in \bw\sbullet(W^*)$, $\beta\in \bw\sbullet(E^*)$,
\begin{equation*}\label{eq:defBere}
\begin{split}
\int^B\alpha\beta&=0, \quad\text{if } \deg \beta<n,\\
\int^B\alpha
e^1\wedge\cdots\wedge e^n&=\frac{(-1)^{n(n+1)/2}}{\pi^{\sfrac{n}{2}}}\,\alpha.
\end{split}
\end{equation*}
More generally, if $o(E)$ is the orientation line of $E$, then $\int^B$ defines a linear map from $\bw\sbullet(W^*)\wotimes\bw\sbullet(E^*)$ into $\bw\sbullet(W^*)\wotimes o(E)$, which is called a Berezin integral.

Let $A\in \End^{\mathrm{anti}}(E)$ be an antisymmetric endomorphism of $E$. We identify $A$ with
\begin{equation}\label{eq:ALam}
\dot{A}=\frac{1}{2}\sum_{1\l i,j\l n}\<e_i,Ae_j\>e^i\wedge e^j\in \bw2(E^*).
\end{equation}
By definition, the Pfaffian $\Pf[A]\in o(E)$ of $A$ is given by
\begin{equation*}\label{eq:PfA}
\Pf[\sfrac{A}{2\pi}]=\int^B\exp(-\sfrac{\dot{A}}{2}).
\end{equation*}
Clearly, $\Pf[A]$ vanishes if $n$ is odd.

\subsection{Mathai-Quillen formalism}\label{sMQ}
Recall that $X$ is a connected closed smooth manifold of dimension $m$. Let
$E$ be a Euclidean vector bundle of rank $n$ on $X$ with a Euclidean
metric $g^{E}$ and a metric connection $\nabla^{E}$. Let
\[
R^{E}=(\nabla^{E})^{2}\in\Omega^{2}(X,\End^\mathrm{anti}(E))
\]
be its curvature. Denote by $o(E)$ the orientation line
bundle of~$E$.
The Euler form of
$(E,\nabla^{E})$ is given by
\begin{equation*}
e(E,\nabla^{E})=\Pf[\sfrac{R^{E}}{2\pi}] \in
\Omega^{n}(X,o(E)).
\end{equation*}
Clearly, $e(E,\nabla^{E})=0$, if $n$ is odd.

Let $\cE$ be the total space of $E$, and let $\pi: \cE\to X$ be the
natural projection. We~will use the formalism of the Berezin integral developed in
Section \ref{sec:Bere} with $W = T\cE$. If~$\omega$ is a smooth section of
$\bw{\sbullet}(T^{*}\cE)\otimes\pi^{*}\bw{\sbullet}(E^{*})$
over $\cE$, then $\int^{B}\omega$ is a smooth section of $\bw{\sbullet}
(T^{*} \cE)
\otimes \pi^{*}o(E)$ over $\cE$.

Let $T^{V}\cE\subset T\cE$ be the vertical subbundle of $T\cE$, and
let $T^{H}\cE\subset T\cE$ be the horizontal subbundle of $T\cE$
with respect to $\nabla^{E}$, so that
\begin{equation}\label{splitting}
T\cE=T^{H}\cE\oplus T^{V}\cE.
\end{equation}
By the identification
$T^{V}\cE\simeq \pi^{*}E$,
the
vertical projection with
respect to the splitting~\eqref{splitting} induces a section
$P^{E}\in C^{\infty}(\cE,T^{*}\cE \otimes \pi^{*}E)$. Using the metric $g^{E}$, we~iden\-tify~$P^{E}$ with $\dot{P}^{E}\in C^{\infty}(\cE,T^{*}\cE \otimes \pi^{*}E^{*})$.
Let $Y\in C^{\infty}(\cE,\pi^{*}E)$ be the tautological section.
Write $\widehat{Y}$ the corresponding section in
$C^{\infty}(\cE,\pi^{*}E^{*})$ induced by $g^{E}$.
Recall that $\dot{R}^{E}\in C^{\infty}(X,\bw{2}(T^{*}X)\otimes \bw{2}(E^{*}))$ is
defined in \eqref{eq:ALam}.

\begin{defin}
For
$T\g 0$, set
\begin{equation*}
A_{T}=\frac{1}{2}\,\pi^{*}\dot{R}^{E}+\sqrt{T}\,\dot{P}^{E}+T|Y|^{2}\in C^{\infty}\big(\cE,\bw{\sbullet}(T^{*}\cE)\otimes \pi^{*}\bw{\sbullet}(E^{*})\big).
\end{equation*}
Let $(\alpha_{T})_{T\g0}$ and $(\beta_{T})_{T>0}$ be families of forms on $\cE$ defined
by
\begin{equation}\label{eqDang1}
\begin{aligned}
&\alpha_{T}=\int^{B}\exp(-A_{T})\in
\Omega^{n}(\cE,\pi^{*}o(E)),\\
&\beta_{T}=\int^{B}\frac{\widehat{Y}}{2\sqrt{T}}\,\exp(-A_{T})\in \Omega^{n-1}(\cE,\pi^{*}o(E)).
\end{aligned}
\end{equation}
\end{defin}
Clearly,
\begin{equation*}
\alpha_{0}=\pi^{*}e(E,\nabla^{E}).
\end{equation*}
Let us recall \cite[Th.\,3.4, 3.5, \& 3.7]{BZ92}.
\begin{thm}
For $T\g 0$, the form $\alpha_{T}$ is closed whose cohomology class
does not depend on $T$.
For $T>0$, $\alpha_{T}$ represents the Thom
class of $E$ so that $\pi_{*}\alpha_{T}=1$, and we have
\begin{equation*}
\frac{\p \alpha_{T}}{\p T}=-d\beta_{T}.
\end{equation*}
\end{thm}

We identify $X$ as a submanifold of $\cE$ by the zero section. The normal bundle
to $X$ in $\cE$ is exactly $E$ and the conormal bundle is $E^{*}$.
Let $\delta_{X}$ be the current on $\cE$ defined by the integration
on $X$. If $\mu$ is a smooth compactly supported form on $\cE$ with
values in $\pi^{*}o(TX)$, then
\begin{equation*}
\int_{\cE}\mu\delta_{X}=\int_{X}\mu.
\end{equation*}
For a current $v$ on $\cE$, denote by $\WF(v)\subset
T^{*}\cE$ its wave front set \cite[\S 8.1]{HormanderBook1}.

\begin{thm}
Let $K\subset
\cE$ be a compact subset of $\cE$. There is $C_{K}>0$ such that for
any $\mu\in
\Omega^{\sbullet}(\cE,\pi^{*}o(TX))$ whose support is contained in $K$
and for any $T\g 1$, we have
\begin{equation}\label{eqabT}
\left| \int_{\cE}\mu(\alpha_{T}-\delta_{X})\right|\l
\frac{C_{K}}{\sqrt{T}}\,\|\mu\|_{C^{1}},\quad\left|\int_{\cE}\mu \beta_{T}\right|\l
\frac{C_{K}}{T^{3/2}}\,\|\mu\|_{C^{1}},
\end{equation}
where $\|{\cdot}\|_{C^1}$ denotes the $C^{1}$-norm.
\end{thm}

In view of \eqref{eqDang1} and \eqref{eqabT},
\begin{equation*}
\psi(E,\nabla^{E})=\int_{0}^{\infty}\beta_{T}dT
\end{equation*}
is a well-defined current of degree $n-1$ on $\cE$ with values in $\pi^{*}o(E)$.

\begin{thm}
The current $\psi(E,\nabla^{E})$ is locally
integrable such that
\begin{equation*}\label{WFphi}
\WF \big(\psi(E,\nabla^{E})\big)\subset E^{*}.
\end{equation*}
The following identity of currents
on $\cE$ holds,
\begin{equation*}
d\psi(E,\nabla^{E})=\pi^{*}e(E,\nabla^{E})-\delta_{X}.
\end{equation*}
\end{thm}

\begin{defin}
The current $\psi(E,\nabla^{E})$ is called the Mathai-Quillen
current.
\end{defin}

\begin{re}
The restriction of $\psi(E,\nabla^{E})$ to the
sphere bundle of $E$
was first constructed in Mathai-Quillen \cite[\S 7]{MathaiQuillen}.
If $E=TX$, this restriction coincides with the transgressed Euler
class defined by Chern \cite{Chern2}.
\end{re}

Assume now $n\l m$. Let $s\in C^{\infty}(X,E)$ be a smooth section of $E$. Set
\begin{equation}\label{eqX0}
X'=\{x\in X:s(x)=0\}.
\end{equation}
Suppose that over $X'$, the differential of $s$ has maximal rank. By transversality,
$X'$ is a smooth submanifold of $X$ of dimension $m-n$. Let $N_{X'/X}$ be the normal bundle to $X'$ in $X$. Using \cite[Th.\,8.2.4]{HormanderBook1}, Bismut and Zhang have shown the following proposition in \cite[Rem.\,3.8]{BZ92}.

\begin{prop}\label{PropBZ38}
The pull-back currents $s^{*}\psi(E,\nabla^{E}),
s^{*}\delta_{X}$ on $X$ are well-defined and satisfy
\begin{equation}\label{eqWFsd}
\WF(s^{*}\psi(E,\nabla^{E}))\subset N_{X'/X}^{*},\quad\WF(s^{*}\delta_{X})\subset N^{*}_{X'/X}.
\end{equation}
The following identity of currents on $X$ holds,
\begin{equation}\label{eqsphi}
d(s^{*}\psi(E,\nabla^{E}))=e(E,\nabla^{E})-s^{*}\delta_{X}.
\end{equation}
\end{prop}

\begin{re}\label{reBZ38}
If $U\in C^{\infty}(X,TX)$ is a vector field on $X$ which has only isolated non degenerated
zeros, \ie in a neighbourhood of a zero $x$ of $U$ there is a
system of coordinates $y=(y_{1},\ldots,y_{m})$ and a matrix
$A$ with $\det
A\neq0$ such that $x$ is represented by $y=0$ and
\begin{equation*}\label{eqU=Ay}
U(y)=Ay+\cO(|y|^{2}).
\end{equation*}
By Proposition \ref{PropBZ38}, $U^{*}\psi(TX,\nabla^{TX})$ is a well
defined current. Moreover, if $\e_{x}=\sgn\det(A)$ is the Poincaré-Hopf
index\footnote{If $U$ is Morse-Smale, then $(-1)^{m}\e_{x}=(-1)^{\ind(x)}.$ See the discussion
after \eqref{eqindxEu} about the sign $(-1)^{m}$.} at $x$, then
$U^{*}\delta_{X}$ is a Radon measure on~$X$ defined for $\mu\in
C^{\infty}(X)$ by
\begin{equation}\label{eqUindd}
\int_{X}\mu \ U^{*}\delta_{X}=\sum_{x:\text{ zero of }U}\e_{x}\,\mu(x).
\end{equation}
\end{re}

\subsection{Ray-Singer metric}\label{sRS}
We use the notation in Section \ref{S1M}. Recall that $(F,\nabla^{F})$ is a flat vector bundle on $X$.
Let $(\Omega^{\sbullet}(X,F),d^{X})$ be the de~Rham complex of smooth
forms on $X$ with values in
$F$. By de~Rham's theory, its cohomology is $H^{\sbullet}(X,F)$.

Take metrics $g^{TX}$ and $g^{F}$ on $TX$ and $F$. Let
$\<,\>_{\bw{\sbullet}(T^{*}X)\otimes F}$ be the induced metric on
$\bw{\sbullet}(T^{*}X)\otimes F$. Let $dv_{X}\in
\Omega^{m}(X,o(TX))$ be the Riemannian volume form on $X$.
For $s_{1},s_{2}\in \Omega^{\sbullet}(X,F)$, set
\begin{equation}\label{eqX}
\<s_{1},s_{2}\>_{\Omega^{\sbullet}(X,F)}=\int_{x\in X}\<s_{1}(x),s_{2}(x)\>_{\bw{\sbullet}(T^{*}X)\otimes
F}dv_{X}.
\end{equation}
Then \eqref{eqX} defines an $L^{2}$-metric on $\Omega^{\sbullet}(X,F)$.

Let $d^{X*}$ be the formal adjoint of $d^{X}$ with respect to the
$L^{2}$-metric $\<\cdot,\cdot\>_{\Omega^{\sbullet}(X,F)}$. Put
\begin{equation*}
\Box^{X}=d^{X}d^{X*}+d^{X*}d^{X}.
\end{equation*}
Then $\Box^{X}$ is a formally self-adjoint second order elliptic differential
operator acting on $\Omega^{\sbullet}(X,F)$. By Hodge theory, we have
\begin{equation}\label{Hodge}
\ker \Box^{X}\simeq H^{\sbullet}(X,F).
\end{equation}
By \eqref{eqlambda} and \eqref{Hodge}, the restriction of the $L^{2}$-metric $\<\cdot,\cdot\>_{\Omega^{\sbullet}(X,F)}$ to $\ker
\Box^{X}$ induces a metric $|\cdot|^{\RS,2}_{\lambda}$ on $\lambda$.

Let $(\ker \Box^X)^{\bot}$ be the orthogonal space to $\ker \Box^X$
in $\Omega^{\sbullet}(X,F)$. Then $\Box^{X}$ acts as an invertible
operator on $(\ker\Box^{X})^{\bot}$. Let $(\Box^{X})^{-1}$ be the
inverse of $\Box^{X}$ acting on $(\ker\Box^X)^{\bot}$. Let
$N^{\sbw{\sbullet}(T^{*}X)}$ be the number operator on
$\bw{\sbullet}(T^{*}X)$, which is multiplication by $p$ on
$\bw{p}(T^{*}X)$. For $s\in \bC$ such that $\Re(s)>\sfrac{m}{2}$, set
\begin{equation*}
\zeta(s)=-\Tr[(-1)^{N^{\sbw{\sbullet}(T^{*}X)}}N^{\sbw{\sbullet}(T^{*}X)}(\Box^{X})^{-s}].
\end{equation*}
By a result of Seeley \cite{Seeley66} or by \cite[Th.\,7.10]{BZ92}, $\zeta(s)$ extends to a meromorphic function
of $s\in \bC$, which is holomorphic at $s=0$.

\begin{defin}
The Ray-Singer metric $\|{\cdot}\|^{\RS,2}_{\lambda}$ on
$\lambda$ is defined by
\begin{equation*}
\|{\cdot}\|^{\RS,2}_{\lambda}=|\cdot|^{\RS,2}_{\lambda}\exp\big(\zeta'(0)\big).
\end{equation*}
\end{defin}

Let $\nabla^{TX}$ be the Levi-Civita connection on $(TX,g^{TX})$. Let
$\psi(TX,\nabla^{TX})$ be the Mathai-Quillen current. By \eqref{eqMS} and by Proposition \ref{PropBZ38}, for any Morse-Smale vector field
$V$, the pull-back $V^{*}\psi(TX,\nabla^{TX})$ is a well-defined current of degree $m-1$ on $X$ with values in $o(TX)$.
Set
\begin{equation*}
\theta(F,g^{F})=\Tr[(g^{F})^{-1}\nabla^{F}g^{F}]\in
\Omega^{1}(X).
\end{equation*}
Then, $\theta(F,g^{F})$ is a closed $1$-form and its cohomology class
$\theta(F)=[\theta(F,g^{F})]\in H^{1}(X)$ does not depend on the
metric $g^F$. Up to a normalization, the class $\theta(F)$ coincides with the first Kamber-Tondeur class \cite{KamberTondeur74}.

\pagebreak[2]
The main result of Bismut-Zhang \cite[Th.\,0.2]{BZ92} is the following.

\begin{thm}\label{BZ}
Suppose that $f$ is a Morse function, whose gradient $\nabla f$ with respect to
some Riemannian metric is Morse-Smale. The following identity holds,
\begin{equation*}\label{eqBZ}
\log\bigl(\sfrac{\|{\cdot}\|^{\RS,2}_{\lambda}}{\|{\cdot}\|^{\rM,2}_{\lambda,-\nabla f}}\bigr)=-\int_{X}\theta(F,g^{F})(\nabla f)^{*}\psi(TX, \nabla ^{TX}).
\end{equation*}
\end{thm}

\subsection{A variation formula for certain characteristic
form}\label{sMa}Let us follow \cite[\S VI.a) and VI.b)]{BZ92}.
Let $(U_{\ell})_{0\l \ell\l 1}$ be a smooth family of vector fields
on $X$, such that each $U_{\ell}$ has only isolated non degenerated
zeros. By Proposition \ref{PropBZ38} and Remark \ref{reBZ38}, the integral
\begin{equation*}
\int_{X}\theta(F,g^{F})U_{\ell}^{*}\psi(TX, \nabla ^{TX})
\end{equation*}
is well-defined. Let us study its variation with respect to $\ell\in [0,1]$.

Let $q: [0,1]\times X\to X$ be the obvious projection. Consider a
smooth section $U\in C^{\infty}([0,1]\times X, q^{*}(TX))$ defined by
\begin{equation*}
U: (\ell,x)\in [0,1]\times X\to U_{\ell}(x)\in T_{x}X.
\end{equation*}
By the consideration after \eqref{eqX0}, the zero set of $U$ is a
manifold of dimension $1$.
Therefore, if $x_{1,0}, \ldots, x_{N,0}$ are the zeros of $U_{0}$,
we can parametrize the zeros of $U_{\ell}$ by
$x_{1,\ell},\ldots,x_{N,\ell}$ such that all the maps $\ell\in
[0,1]\to x_{i,\ell}\in X$ are
smooth. Also, the Poincaré-Hopf index of $U_{\ell}$ at
$x_{i,\ell}$ does not depend on $\ell$ and will be denoted by
$\e_{i}\in \{\pm1\}$. The following proposition is a generalization of \cite[Prop.\,6.4]{BZ92}. Since we will use this proposition several times in
Section \ref{s33}, let us give a detailed proof.

\begin{prop}\label{propBZVVa}
The following identity holds:
\begin{equation}\label{eqDDMA}
\int_{X}\theta(F,g^{F})
\Big(U_{1}^{*}\psi(TX,\nabla^{TX})-
U_{0}^{*}\psi(TX,\nabla^{TX})\Big)
=\sum_{i=1}^{N}\e_{i}\int_{0}^{1}\theta(F,g^{F})(\dot{x}_{i,\ell})d\ell.
\end{equation}
\end{prop}
\begin{proof}
Let us follow the proof of \cite[Prop.\,6.1, Prop.\,6.4]{BZ92}. Equip the pull-back vector bundle
$q^{*}(TX)$ over $[0,1]\times X$ with the pull-back metric and the pullback metric
connection $\nabla^{q^{*}(TX)}$. Let
$\psi(q^{*}(TX),\nabla^{q^{*}(TX)})$ be the corresponding
Mathai-Quillen current. By Proposition \ref{PropBZ38},
$U^{*}\psi(q^{*}(TX),\nabla^{q^{*}(TX)})$ and
$U^{*}\delta_{[0,1]\times X}$ are well-defined currents such that
\begin{equation}\label{eqdphi01}
d^{[0,1]\times X} (
U^{*}\psi(q^{*}(TX),\nabla^{q^{*}(TX)}))=e(q^{*}(TX),\nabla^{q^{*}(TX)})-U^{*}\delta_{[0,1]\times X}.
\end{equation}
By our construction,
\begin{equation}\label{eqdphi02}
e(q^{*}(TX),\nabla^{q^{*}(TX)})=q^{*}e(TX,\nabla^{TX}).
\end{equation}
Since $\theta(F,g^{F})$ is a closed $1$-form on $X$, by
\eqref{eqdphi01} and \eqref{eqdphi02}, we have
\begin{equation*}
d^{[0,1]\times X} (q^{*}\theta(F,g^{F}) \wedge
U^{*}\psi(q^{*}(TX),\nabla^{q^{*}(TX)}))
=q^{*}\theta(F,g^{F}) \wedge U^{*}\delta_{[0,1]\times X}.
\end{equation*}
Integrating the above formula over $[0,1]\times X$, by the Stokes
formula, we get \eqref{eqDDMA}.
\end{proof}

\subsection{Proof of the main result}\label{s33}

We restate our main
result Theorem \ref{thm1}, which is an extension of Theorem \ref{BZ}.
\begin{thm}\label{Thm1}
Suppose that $V$ is a Morse-Smale vector field. The following identity holds,
\begin{equation*}
\log\bigl(\sfrac{\|{\cdot}\|^{\RS,2}_{\lambda}}{\|{\cdot}\|^{\rM,2}_{\lambda,V}}\bigr)=-\int_{X}\theta(F,g^{F})(-V)^{*}\psi(TX,\nabla^{TX}).
\end{equation*}
\end{thm}
\begin{proof}
Take $(V_{\ell})_{0\l \ell\l 1}$ as in Proposition
\ref{propV01}.
By Remark \ref{reDMA}, we have
\begin{equation}\label{MVMV1}
\|{\cdot}\|^{\rM,2}_{\lambda,V}=\|{\cdot}\|^{\rM,2}_{\lambda,V_{1}}.
\end{equation}
Since the critical elements of $V$ and $V_{1}$ coincide, the fixed points of $V_{\ell}$ form smooth loops on $X$. By Remark \ref{reVl}, we can assume that the fixed points of $V_{\ell}$ are in a small neighbourhood of
the fixed points set of $V$. In particular, the above loops are
contractible. By Proposition \ref{propBZVVa} and by the closedness of $\theta(F,g^{F})$, we have
\begin{equation}\label{XVV1}
\int_{X}\theta(F,g^{F})(-V)^{*}\psi(TX,
\nabla^{TX})=\int_{X}\theta(F,g^{F})(-V_{1})^{*}\psi(TX, \nabla^{TX}).
\end{equation}
By \eqref{MVMV1} and \eqref{XVV1}, it is enough to show our theorem
for the Morse-Smale vector field $V$ whose critical elements are all
of standard forms.

Take $f$ as in Proposition
\ref{propfV}.
By Proposition \ref{propMM} and Theorem \ref{BZ}, we have
\begin{multline}\label{LLa}
\log\bigl(\sfrac{\|{\cdot}\|^{\RS,2}_{\lambda}}{\|{\cdot}\|^{\rM,2}_{\lambda,V}}\bigr)=-\int_{X}\theta(F,g^{F})(\nabla f)^{*}\psi(TX, \nabla ^{TX})\\
+\sum_{\gamma\in
B}(-1)^{\ind(\gamma)}\log
\left\|{\det\,}\tau(a_{\gamma}')\right\|^{2}_{\det
F_{x'_{\gamma}}\otimes (\det
F_{x_{\gamma}})^{-1}}.
\end{multline}
By \eqref{LLa}, it remains to show
\begin{multline}\label{eqaaa}
\int_{X}\theta(F,g^{F})(\nabla
f)^{*}\psi(TX,
\nabla^{TX})-\int_{X}\theta(F,g^{F})(-V)^{*}\psi(TX,\nabla^{TX})\\
=\sum_{\gamma\in
B}(-1)^{\ind(\gamma)}\log
\left\|{\det\,}\tau(a_{\gamma}')\right\|^{2}_{\det
F_{x'_{\gamma}}\otimes (\det
F_{x_{\gamma}})^{-1}}.
\end{multline}
We assume that all the closed orbits are in standard form
of Case \ref{Case:1}. Cases \ref{Case:2}--\ref{Case:4} can be dealt similarly.

Following \cite[\S IV.c)]{BZ92},
choose a smooth triangulation $K$ of $X$ such that $A\cap
K^{m-1}=\varnothing$, and such that on $\ol{U}_{\gamma}\simeq \bbS^{1}\times
\ol{\bbD}^{m-1}$ the
triangulation is given by the $m$-simplex
$\sigma^{m}_{\gamma}=(\bbS^{1}-\{1/4\})\times
\ol{\bbD}^{m-1}$ and $(m-1)$-simplex
$\sigma^{m-1}_{\gamma}=\{1/4\}\times \ol{\bbD}^{m-1}$.

On each simplex $\sigma\in K^{m}\setminus K^{m-1}$ of maximal degree, choose a
primitive $W_{0,\sigma}\in C^{\infty}(\sigma)$ of $\theta(F,g^{F})$,
such that on $\sigma$ we have
\begin{equation*}
dW_{0,\sigma}=\theta(F,g^{F}).
\end{equation*}
Let $W_{0}$ be the locally integrable current on $X$, such that
for each $\sigma\in K^{m}\setminus K^{m-1}$ we have
\begin{equation*}
W_{0}|_{\sigma}=W_{0,\sigma}.
\end{equation*}
By our construction of $K$, for $\gamma\in B$, the two points
$x_{\gamma},x_{\gamma}'$, and the integral curve $a_{\gamma}'$ are in the same simplex $\sigma^{m}_{\gamma}\in
K^{m}$, so that
\begin{equation}\label{W0W0}
W_{0}(x_{\gamma}')-W_{0}(x_{\gamma})=\log
\left\|{\det\,}\tau(a_{\gamma}')\right\|^{2}_{\det
F_{x'_{\gamma}}\otimes (\det
F_{x_{\gamma}})^{-1}}.
\end{equation}

Set
\begin{equation}\label{eqW1}
W_{1}=\theta(F,g^{F})-dW_{0}.
\end{equation}
Then $W_{1}$ is a closed
current of degree $1$ on $X$ such that $\Supp(W_{1})\subset
K^{m-1}$. By \eqref{eqWFsd} and by $A\cap K^{m-1}=\varnothing$, $(-V)^{*}\psi(TX,\nabla^{TX})$ is
smooth in the neighbourhood of the support of $W_{1}$, so that
\begin{equation*}
W_{1}\wedge (-V)^{*}\psi(TX,\nabla^{TX})
\end{equation*}
is a well-defined current on $X$. By \eqref{eqsphi}, \eqref{eqUindd}, and \eqref{eqW1}, we have
\begin{multline}\label{eqVtoW1}
-\int_{X}\theta(F,g^{F}) (-V)^{*}\psi(TX,\nabla^{TX})\\
=\int_{X}W_{0}\ e(TX,\nabla^{TX})
-\sum_{x\in A}(-1)^{\ind(x)}W_{0}(x)
-\int_{X}W_{1}\wedge (-V)^{*}\psi(TX,\nabla^{TX}).
\end{multline}
Similar when $-V$ is replaced by $\nabla f$, we get
\begin{multline}\label{eqnftoW1}
-\int_{X}\theta(F,g^{F}) (\nabla f)^{*}\psi(TX,\nabla^{TX})\\
=\int_{X}W_{0}\ e(TX,\nabla^{TX})
-\sum_{x\in A}(-1)^{\ind(x)}W_{0}(x)\\
+\sum_{\gamma\in
B}(-1)^{\ind(\gamma)}(W_{0}(x_{\gamma})-W_{0}(x_{\gamma}')) -\int_{X}W_{1}\wedge (-V)^{*}\psi(TX,\nabla^{TX}).
\end{multline}

By \eqref{W0W0}, \eqref{eqVtoW1}, and \eqref{eqnftoW1}, we see that \eqref{eqaaa} is equivalent to
\begin{equation}\label{eqLast}
\int_{X}W_{1} \wedge (\nabla
f)^{*}\psi(TX,\nabla^{TX})=\int_{X}W_{1}\wedge
(-V)^{*}\psi(TX,\nabla^{TX}).
\end{equation}
By \eqref{eqdf=V}, on any simplex in
$K^{m-1}$ other than $\sigma_{\gamma}^{m-1}$, we have $\nabla f=-V$.
By Remark \ref{reK}, near $\sigma_{\gamma}^{m-1}$, $\nabla
f$ and $-V$ can be connected by a family of vector fields without
zero. Using the fact that $\Supp(W_{1})\subset K^{m-1}$,
and by a version of Proposition \ref{propBZVVa} where
$\theta(F,g^{F})$ is replaced by the closed current $W_{1}$, we get \eqref{eqLast}.
The proof of our theorem is completed.
\end{proof}

